% EXTRACTED from v1.1 manuscript (tag v1.1-verified): Part II, ch. "Natural
% Deduction and Sequents" (ND_CPL rule displays), presented in the contexted
% form of the "Rule Sheets" chapter (sec:rule-sheets-01), which the manuscript
% states is literally the active calculus with contexts made explicit.
% The calculus is the tree calculus of the source derivation definition
% (restoration/inherited-02193-02340.tex:1-30) with derivability as in
% inherited-02193-02340.tex:80-86 (open assumptions all belong to Gamma).
% Each contexted display is read as one rule application on witnessing
% derivations, so the discharge side conditions speak about those witnesses.
% Discharge follows inherited-02193-02340.tex:14-21 exactly: an application
% discharges exactly the open occurrences of its bracketed formula marked with
% its own label (possibly some but not all of them), every carried label
% discharges at least one occurrence, and the carried label need not be fresh;
% freshness is required only for graft copies (8-12) and for the merging
% relabeling (21-27). There is no repetition rule (27-29).
\calculussheet{sheet:nd-cpl}{Classical Propositional Natural Deduction ($\mathsf{ND}_{\mathsf{CPL}}$)}{Language: $\mathsf{Fm}_{\mathsf{CPL}}$ with primitives $\bot$, $\neg$, $\land$, $\lor$, $\to$; derivations: finite rooted trees of formula occurrences whose leaves are assumption occurrences and whose other nodes are applications of the twelve rules displayed after the assumption display; discharge labels are metalinguistic natural numbers: a leaf may carry a label, an application of $\to I$, $\neg I$ or RAA carries one label, an application of $\lor E$ carries two, and an assumption occurrence is open iff no rule application on its path to the root discharges it; sequent shape: single-conclusion judgments $\Gamma\vdash\varphi$ with $\Gamma$ a finite set of formulas, where $\Gamma\vdash\varphi$ holds iff some derivation has root $\varphi$ and every open assumption occurrence in it is an occurrence of a member of $\Gamma$; members of $\Gamma$ need not occur, and since a derivation has finitely many leaves, derivability from an arbitrary set of formulas is derivability from a finite subset of it; reading of the displays: each display is one step on witnessing derivations, not a closure condition on the relation alone: every premise of an instance stands for a chosen derivation witnessing it, one application of the named rule is placed beneath the roots of these derivations, and the result witnesses the conclusion; the assumption display is the one-node derivation $\varphi$, with or without a label, whose single open occurrence is of a member of $\Gamma\cup\{\varphi\}$, and there is no separate weakening rule; discharge: for a premise displayed as $\Delta\cup\{\varphi\}\vdash\chi$ of a discharging rule, the application carries a label $i$ and discharges exactly the open occurrences of $\varphi$ marked $i$ in the derivation witnessing that premise; at least one such occurrence must exist, and they may be some but not all of the open occurrences of $\varphi$ there; the carried label need not be fresh, because what an application discharges is fixed by the marks on the open occurrences of its own bracketed formula above its own premise, so an open occurrence of another formula marked $i$, and an open occurrence of $\varphi$ with another label or none, stays open; the result witnesses the displayed conclusion only if every open occurrence that remains in it is of a member of the conclusion context, so an open occurrence of $\varphi$ left unmarked is allowed exactly when $\varphi$ belongs to that context; to discharge every open occurrence of $\varphi$ in the witness, first mark them all with one common fresh label by convention~(ii) and let the application carry that label; vacuous discharge is not licensed: a witness with no open occurrence of $\varphi$ marked $i$ gives no instance of the display with label $i$; if the witness has no open occurrence of $\varphi$ at all, it already shows $\Delta\vdash\chi$, and otherwise convention~(ii) supplies a label marking one; label conventions, neither of which is an inference step, and which are the only places where a label must be fresh: (i)~applying one injective map to all assumption labels of a derivation and to the matching labels on its rule applications preserves the derivation and the occurrences each application discharges, and when copies of derivations are grafted onto open assumptions (Lemma~\ref{lem:nd-substitution}) each copy is first renamed into a label set disjoint from every label of the receiving derivation and of every other copy; (ii)~merging relabeling: marking any set of open occurrences of one formula with a common label fresh from every label of the derivation preserves the derivation and the occurrences each application discharges, because no application carries that label, and this is how a newly added final inference discharges several open occurrences at once (as in the proof of Theorem~\ref{thm:prop-completeness}, which uses $\bot E$ when no open occurrence of the negated goal remains and otherwise this relabeling followed by RAA); there is no repetition rule: a repeated formula is the node already present; negation: primitive, with its own rules alongside primitive $\bot$, and RAA supplying classical strength; equality: none (propositional language).}
\hbrule{rule:nd-cpl-assumption}{assumption}{}{\Gamma\cup\{\varphi\}\vdash\varphi}{\none}
\hbrule{rule:nd-cpl-and-i}{$\land I$}{\Gamma\vdash\varphi\qquad\Delta\vdash\psi}{\Gamma\cup\Delta\vdash\varphi\land\psi}{\none}
\hbrule{rule:nd-cpl-and-e1}{$\land E_1$}{\Gamma\vdash\varphi\land\psi}{\Gamma\vdash\varphi}{\none}
\hbrule{rule:nd-cpl-and-e2}{$\land E_2$}{\Gamma\vdash\varphi\land\psi}{\Gamma\vdash\psi}{\none}
\hbrule{rule:nd-cpl-or-i1}{$\lor I_1$}{\Gamma\vdash\varphi}{\Gamma\vdash\varphi\lor\psi}{\none}
\hbrule{rule:nd-cpl-or-i2}{$\lor I_2$}{\Gamma\vdash\psi}{\Gamma\vdash\varphi\lor\psi}{\none}
\hbrule{rule:nd-cpl-or-e}{$\lor E$}{\Gamma\vdash\varphi\lor\psi\qquad\Delta\cup\{\varphi\}\vdash\chi\qquad\Sigma\cup\{\psi\}\vdash\chi}{\Gamma\cup\Delta\cup\Sigma\vdash\chi}{discharge: the application carries two labels $i,j$, as displayed, neither of which need be fresh; it discharges exactly the open occurrences of $\varphi$ marked $i$ in the derivation witnessing the second premise and exactly the open occurrences of $\psi$ marked $j$ in the one witnessing the third, and each of these two sets is nonempty; an open occurrence of $\varphi$ (of $\psi$) left unmarked there stays open, so the instance requires $\varphi$ (respectively $\psi$) to belong to $\Gamma\cup\Delta\cup\Sigma$ whenever one is left}
\hbrule{rule:nd-cpl-imp-i}{$\to I$}{\Gamma\cup\{\varphi\}\vdash\psi}{\Gamma\vdash\varphi\to\psi}{discharge: the application carries a label $i$, which need not be fresh, and discharges exactly the open occurrences of $\varphi$ marked $i$ in the derivation witnessing the premise, of which there is at least one; an open occurrence of $\varphi$ left unmarked stays open, so the instance requires $\varphi\in\Gamma$ whenever one is left}
\hbrule{rule:nd-cpl-imp-e}{$\to E$}{\Gamma\vdash\varphi\to\psi\qquad\Delta\vdash\varphi}{\Gamma\cup\Delta\vdash\psi}{\none}
\hbrule{rule:nd-cpl-neg-i}{$\neg I$}{\Gamma\cup\{\varphi\}\vdash\bot}{\Gamma\vdash\neg\varphi}{discharge: the application carries a label $i$, which need not be fresh, and discharges exactly the open occurrences of $\varphi$ marked $i$ in the derivation witnessing the premise, of which there is at least one; an open occurrence of $\varphi$ left unmarked stays open, so the instance requires $\varphi\in\Gamma$ whenever one is left}
\hbrule{rule:nd-cpl-neg-e}{$\neg E$}{\Gamma\vdash\varphi\qquad\Delta\vdash\neg\varphi}{\Gamma\cup\Delta\vdash\bot}{\none}
\hbrule{rule:nd-cpl-bot-e}{$\bot E$}{\Gamma\vdash\bot}{\Gamma\vdash\varphi}{\none}
\hbrule{rule:nd-cpl-raa}{RAA}{\Gamma\cup\{\neg\varphi\}\vdash\bot}{\Gamma\vdash\varphi}{discharge: the application carries a label $i$, which need not be fresh, and discharges exactly the open occurrences of $\neg\varphi$ marked $i$ in the derivation witnessing the premise, of which there is at least one; an open occurrence of $\neg\varphi$ left unmarked stays open, so the instance requires $\neg\varphi\in\Gamma$ whenever one is left; this is the explicitly classical rule, and omitting it gives the primitive-negation presentation of the intuitionistic fragment}
